\chapter{Introduction}

\section{Problem Statement}

Fine-grained artifact recognition in museum environments is difficult due to
background clutter, low-light capture, visual overlap between classes, and open-set
uncertainty. At the \GEM{} (GEM) the difficulty is acute: the target set contains
\textbf{51 catalogued artifacts}, many of which are near-identical Ramesside statues
that differ only in subtle pose, headdress, or accompanying standards. Single-model
solutions trained on idealized imagery often produce unstable predictions in field
conditions.

In practical terms, this means the system can appear confident when it should be
cautious, and uncertain when the visitor expects immediate clarity. A tourist
standing in front of two visually similar statues does not care about model
internals; they care about whether the app gives a trustworthy answer quickly, and
in their own language. This mismatch between academic benchmark success and
real-world reliability is the true problem addressed by this thesis.

\section{Motivation and Significance}

\TimeLens{} is motivated by practical museum needs: reliable visitor guidance,
catalogue support, and safe AI--human interaction. The project targets a
production-ready path where on-device recognition is strong enough to drive a
downstream knowledge experience without sending every camera frame to the cloud.

The significance is both technical and cultural. Technically, the work advances a
\emph{label-quality-driven} recipe for fine-grained visual recognition under
difficult conditions, and a \emph{grounded} bilingual question-answering layer that
resists fabrication. Culturally, it contributes to accessible heritage
interpretation by helping visitors understand artifacts in context, in real time,
and in natural language. The thesis treats this dual objective seriously: high model
quality without user trust is insufficient, and polished UX without reliable identity
and grounding is unsafe.

\begin{insightbox}[The TimeLens Vision]
We don't just recognize artifacts; we unlock their stories. By pairing a compact
\textbf{on-device detector} with a \textbf{retrieval-grounded Gemma guide}, we turn
every smartphone into a bilingual curatorial companion.
\end{insightbox}

\section{Research Gaps, Questions, and Objectives}

\subsection{Research Gaps}
\begin{itemize}[leftmargin=1.5em]
  \item Generic detectors trained on web imagery are weak on fine-grained,
        domain-specific artifacts such as near-identical Ramesside statues.
  \item On datasets built from video frames, naïve random train/validation splits
        leak near-duplicate frames across the split and inflate reported accuracy.
  \item Public-facing museum AI guides hallucinate historical facts when generation
        is not grounded in curated data, and frequently handle Arabic poorly.
\end{itemize}

\subsection{Research Questions}
\begin{itemize}[leftmargin=1.5em]
  \item \textbf{RQ1.} Can label-quality iteration (auto-annotation $\rightarrow$
        cleaning $\rightarrow$ hand-annotation) yield a compact on-device detector
        that reliably distinguishes 51 fine-grained Egyptian artifacts?
  \item \textbf{RQ2.} Can a retrieval-grounded, locally-served language model deliver
        factual, bilingual (English/Arabic) museum answers at interactive latency?
  \item \textbf{RQ3.} Can the combined system run on commodity phones and a single
        small GPU, making it suitable for real museum deployment?
\end{itemize}

\subsection{Objectives}
\begin{itemize}[leftmargin=1.5em]
  \item Build and evaluate a compact, leakage-aware on-device artifact detector
        (YOLOv8n) exported to TensorFlow Lite.
  \item Build a grounded bilingual RAG guide (ChromaDB retrieval + Gemma 4 E2B
        generation) served behind a FastAPI interface.
  \item Integrate both in a Flutter application and document the market relevance of
        AR plus deep learning for heritage technology.
\end{itemize}

\section{Proposed Solution Overview}

\TimeLens{} splits cleanly into an on-device perception layer and a grounded
knowledge layer. On the phone, a compact \textbf{YOLOv8n} detector---exported to
TensorFlow Lite with end-to-end non-maximum suppression---runs on the live camera
feed and classifies the framed exhibit as one of 51 catalogued artifacts, drawing a
tracking box in real time. Because the detector was trained on clean, hand-verified
labels, its confidence is well calibrated: real artifacts score high while off-target
scenes fall below a rejection threshold, so the app can \textbf{trust the on-device
label directly} without a per-frame network round-trip. When the visitor taps a
confident detection, the app opens the artifact's information and a conversational
guide; views that fall below the per-class confidence threshold produce no overlay at
all, so the app simply keeps scanning rather than guessing.

The conversational guide is a bilingual Retrieval-Augmented Generation (RAG)
service. Rather than relying on a language model's parametric memory, it retrieves
the most relevant entries from a curated 108-record ChromaDB knowledge base and
passes them to Gemma 4 E2B---served locally through Ollama behind a FastAPI
endpoint---as grounding context. This keeps every answer anchored to verified museum
data in the visitor's language (English or Arabic) and structurally limits invented
historical claims. A backend re-verification of the on-device label is supported by
the architecture but is not required in the deployed path, where the calibrated
detector is the source of truth.

\section{Thesis Organization}

The remainder of the thesis is organized as follows. Chapter~2 reviews related work
in on-device detection, foundation-model auto-annotation, and retrieval-augmented
generation. Chapter~3 presents the analysis and requirements. Chapter~4 describes the
deployed system architecture. Chapter~5 is the on-device detection study (the
data-quality-driven YOLOv8n contribution). Chapter~6 covers the bilingual RAG guide
and its deployment. Chapter~7 is the market analysis of AR and AI guide systems for
Egyptian cultural heritage. Chapter~8 concludes and outlines future work. Appendices
cover ethics, reproducibility, and a generative-AI disclosure.
